\section{Spectral theory in finite dimensional normed spaces}

1 Only do B. Notice that B corresponds to complex numbers. As in if we take all matrices of this type as a ring, you get exactly the complex numbers.
2 This will be seen to be true very generally
3 This will also seen to be true very generally 
4 This too. Remember the discussion of c-numbers (complex numbers) and q-numbers (operators/matrices), where we had the dictionary that
real numbers <-> hermitian
imaginary numbers <-> skew-hermitian 
on unit circle <-> unitary
Problem 2,3,4 show that the eigenvalues behave entirely well according to this dictionary 
5 This is also true in general, notice that the spectrum in general is not just the same as the set of eigenvalues as it is here 
6 Very good intuition for why the trace/determinant is useful, it is the product and sum of eigenvalues.
7
9
10 These three exercises show that "standard" operations on matrices carry over to eigenvalues. Inversing a matrix, inverses the eigenvalues etc. This will be seen to hold very very generally, both for operators, but also for more general functions than polynomials. For example, if A is a matrix, then the eigenvalues of $e^A$ will be exactly of the form $e^lambda$. But of course we'd need to define what the exponential/logarithm/sqrt/... of a matrix is. This is not easy at all.
11 Notice that matrices A and $C^-1$ A C are denoted as similar matrices. This is an equivalence relation. Geometrically it means that these matrices come from the same operator but in a different basis. It is very good that two similar matrices share the same eigenvalues
12 This is the easiest example of a matrix that is not diagonalizable. Functional analysis will not really be able to handle these matrices. It will be more about rather well behaved matrices (hermitian, unitary, ...)

\begin{exercise}{1}
fill
\end{exercise}
\begin{proof}
fill
\end{proof} 

\begin{exercise}{2}
fill
\end{exercise}
\begin{proof}
fill
\end{proof} 

\begin{exercise}{3}
fill
\end{exercise}
\begin{proof}
fill
\end{proof} 

\begin{exercise}{4}
fill
\end{exercise}
\begin{proof}
fill
\end{proof} 

\begin{exercise}{5}
fill
\end{exercise}
\begin{proof}
fill
\end{proof} 

\begin{exercise}{6}
fill
\end{exercise}
\begin{proof}
fill
\end{proof} 

\begin{exercise}{7}
fill
\end{exercise}
\begin{proof}
fill
\end{proof} 

\begin{exercise}{9}
fill
\end{exercise}
\begin{proof}
fill
\end{proof} 

\begin{exercise}{10}
fill
\end{exercise}
\begin{proof}
fill
\end{proof} 

\begin{exercise}{11}
fill
\end{exercise}
\begin{proof}
fill
\end{proof} 

\begin{exercise}{12}
fill
\end{exercise}
\begin{proof}
fill
\end{proof} 
